% DEIXIS evidence report, IEEEtran export for XeLaTeX.
% Compile with: latexmk -xelatex report-synthetic-report-draft.tex
% Needs XeLaTeX, BibTeX and the IEEEtran class and IEEEtran.bst style; they are not included.
% Export notes: none

\documentclass[journal]{IEEEtran}
\usepackage{fontspec}
\usepackage{amsmath}
\usepackage{amssymb}
\usepackage{bm}
\usepackage{xcolor}
\usepackage{cite}
\usepackage{url}
\usepackage{booktabs}
\usepackage{longtable}
\usepackage{array}

\begin{document}

\renewcommand{\abstractname}{Özet}
\renewcommand{\IEEEkeywordsname}{Dizin Terimleri}
\renewcommand{\refname}{Kaynaklar}

\title{SYNTHETIC report \& \ensuremath{\alpha}}

\maketitle

\noindent Kanıt raporu \ensuremath{\cdot} taslak

\noindent TASLAK: 3 bölüm doğrulanmadı.

\noindent\textit{Özet: Bu bölüm doğrulanmadı ve yeniden yazılmalı.}

\begin{abstract}

SYNTHETIC abstract.

\end{abstract}

\noindent\textit{Dizin Terimleri: Bu bölüm doğrulanmadı ve yeniden yazılmalı.}

\begin{IEEEkeywords}

SYNTHETIC terms

\end{IEEEkeywords}

\section*{I. Giriş}

SYNTHETIC claim. \cite{Synth26}

\section*{III. Arka Plan ve Sınıflandırma}

\noindent\textit{Bu bölüm doğrulanmadı ve yeniden yazılmalı.}

Bu bölüm için metin yazılmadı.

\section*{Rapor notları}

DEIXIS ile üretildi.

Atıf çapaları ilgili pasajlarda veya hücrelerde bulundu. Bu rapor için model veya kişi incelemesi kaydedilmedi; kod, her pasajın ilgili iddiayı destekleyip desteklemediğini denetlemedi.

Bu çıktı raporun metnini ve kaynakçasını taşır. Alıntı pasajları, tablo hücresi kanıtları ve konumlanmış atıf çapaları DEIXIS'te kalır.

\bibliographystyle{IEEEtran}

\bibliography{report-synthetic-report-draft}

\end{document}
